\documentclass[11pt]{article}
\usepackage[T1]{fontenc}
\usepackage{lmodern,microtype,amsmath,amssymb,booktabs,array,enumitem}
\usepackage[margin=0.9in]{geometry}
\usepackage[hidelinks]{hyperref}
\hypersetup{pdftitle={Specialist review packet: chain-cycle separation for leximin protection},pdfsubject={Proof, priority and significance review request}}
\setlist{itemsep=4pt,topsep=5pt}
\setlength{\parindent}{0pt}
\setlength{\parskip}{6pt}
\begin{document}
\begin{center}
{\Large\bfseries Specialist review request}\\[5pt]
{\large A chain--cycle separation for leximin protection}\\[3pt]
September 6, 2026\quad\textbullet\quad Two-page review brief
\end{center}

\textbf{Requested assessment.} Please assess mathematical correctness,
historical priority, and significance separately. The manuscript has an
internal proof audit, exact certificates and a Claude snapshot review;
these do not substitute for a domain specialist's assessment.

\textbf{Exact model.} Three sectors of equal finite size; only complete
triples produce value. Every triple is available. Coalition worth is
maximum \emph{integral} matching value. The core includes every labeled
payoff, nonnegativity, efficiency and all triple covers. Let $x$ be its
\emph{full ascending sorted leximin} point, and
$H_i=\max\{y_i:y\text{ is in the same entire core}\}$.

\textbf{The separation.}
\begin{enumerate}
\item \textbf{Strict chains:} if $w_{ijk}=P_{ij}+Q_{jk}$ with $P,Q\geq0$,
then $x_i\geq H_i/3$ for every player, including zero-ideal coordinates.
The constant is attained by the three-player simplex.
\item \textbf{Cyclic pair interactions:} for each $q\geq1$, a market with
$9q+3$ players and pair entries in $\{0,1,2\}$ has true worth $6q+1$,
a strictly positive core allocation, and a target with
$x_i=1/[3(q+1)]$ and $H_i=1$.
Thus the best class-wide protection fractions are $1/3$ and $0$.
\end{enumerate}

\textbf{Two proof mechanisms.} The chain core is a potential system with
payoff arcs between two layers and one anchor. A classical slack-balancing
cut argument produces a simple cycle for every player; each such cycle
contains exactly three payoff arcs. For the cyclic family, an optimal
perfect matching is worth $6q+1$, while a disjoint partial matching is
worth $6q$ and leaves $3q+3$ players. Their total core payoff is at most
one. An explicit cover equalizes them at $1/[3(q+1)]$; another attains
the target's ideal of one. Both covers are verified symbolically for
every mixed triple and every $q$.

\textbf{Priority claim under consideration.} The candidate contribution
is the exact constrained realization and the resulting separation. We
credit the strict assignment models, coordinate-ideal protection,
selected-arc slack balancing and LP/matching certificates to prior work.
An earlier flow-path example already gives vanishing protection in a
broader class. We do not assert unconditional novelty, a new general
method, optimal finite-size ratios, or a field-level breakthrough.

\textbf{Reading order.} Manuscript Theorem 2.2; Section 3 for the chain
certificate; Section 4 for the cyclic construction. The source comparison
and the replay instructions accompany the manuscript.

\newpage
\begin{center}{\Large\bfseries Priority comparison and review questions}\end{center}

\textbf{Original-source status.}
\begin{itemize}
\item \textbf{Mart\'inez-de-Alb\'eniz--Rafels--Ybern (2023), SSRN4369524:}
the complete 13-page preprint was recovered. Definitions and Theorems
1--4 were compared. Its general concatenation representation retains
outside-strip inequalities; simpler descriptions need additional
conditions. It does not state the leximin ratio theorem or the growing
bounded-pair family. It is distinct from the 2020 assortative paper.
\item \textbf{Quint (1991), GEB 3:487--503:} publisher abstract and author
records read; full original remains behind access controls.
\item \textbf{Lucas (1995), Duke Math. J. 81(1):55--65:} citation now
confirmed on the primary publisher page; full original requires access.
\item \textbf{Stuart (1997), MSS 34(1):21--27:} original introduction and
the displayed nonemptiness proposition are visible in the publisher
preview. Full definitions/proof remain unavailable. The chain formula
is verified in later primary author texts.
\end{itemize}

\textbf{Questions that would change the assessment.}
\begin{enumerate}
\item Is there an invalid implication in the full-core representation,
the selected-slack perturbation, or the three-payoff-arc count?
\item Does the cyclic proof omit any mixed-triple constraint, substitute
fractional for integral worth, or confuse a floor witness with full
leximin? Are the first-face pinning and ideal attainment sufficient?
\item Does an earlier theorem or construction already imply this exact
separation? Please identify the original location and the deduction,
including preservation of labeled payoffs, ideals and the sorted
selector. An equivalent construction can defeat priority without using
the same terminology.
\item If the result is correct and new at this scope, is the separation
economically or mathematically substantial enough for a standalone
paper? Which interpretation or missing comparison is most important?
\end{enumerate}

\textbf{Reproduction.} From the supplement root, run
\texttt{bash build.sh} for standard-library exact checks.
With \texttt{z3-solver==5.1.0.0}, run
\texttt{bash build.sh --smt}. The optional mode adds three raw-core
exclusions for the cyclic target and full-leximin/ideal computations
with per-player cycles for three selected chain markets. The all-$q$
result is a conventional proof, not a finite computational extrapolation.

\textbf{Useful response format.} For each objection, give its location,
the missing hypothesis or contrary example, and whether it affects
correctness, priority, significance or exposition. A bibliographic lead
to any inaccessible original would also be valuable.

\vfill
\small This packet has been prepared for review; it has not been sent.
No external specialist endorsement is recorded.
\end{document}
